% =========================================================================
%  AIinRT 2027 — Abstract submission template (LaTeX)
%  Artificial Intelligence in Radiation Therapy
%  25–26 March 2027 · Princess Máxima Center, Utrecht · https://AIinRT.org
%
%  REQUIREMENTS (delete the framed instructions box before submission):
%   • Main text (Introduction–Conclusion): max. 500 words
%   • One A4 page total, including display items and references
%   • Max. 2 display items (figures and/or tables combined), each captioned;
%     captions and references are NOT counted in the 500 words
%   • Do not change fonts, margins, or spacing; do not add pages
%   • DOUBLE-BLIND review: no author names, affiliations, acknowledgements,
%     or identifying info anywhere in this document; cite own prior work in
%     the third person; author details go in the submission system only
%   • AI-use disclosure is mandatory ("None" is acceptable)
%   • Submit as PDF via AIinRT.org (deadline: see website)
%   • Accepted abstracts: proffered papers, 9 min + 3 min discussion
%
%  Word-count check (TeX Live):  texcount -inc -sum main.tex
% =========================================================================
\documentclass[10pt,a4paper]{article}

% ---- Layout (do not modify) ---------------------------------------------
\usepackage[a4paper,top=20mm,bottom=18mm,left=20mm,right=20mm]{geometry}
\usepackage[T1]{fontenc}
\usepackage[utf8]{inputenc}
\usepackage{helvet}\renewcommand{\familydefault}{\sfdefault} % Calibri-like sans
\usepackage{graphicx}
\usepackage{booktabs}
\usepackage{caption}
\captionsetup{font=small,labelfont=bf,skip=4pt}
\usepackage[dvipsnames]{xcolor}
\definecolor{aiNavy}{HTML}{1F3864}
\definecolor{aiAccent}{HTML}{2E74B5}
\definecolor{aiGrey}{HTML}{595959}
\usepackage{titlesec}
\titleformat{\section}{\normalsize\bfseries\color{aiNavy}}{}{0pt}{}
\titlespacing{\section}{0pt}{6pt}{1pt}
\usepackage[hidelinks]{hyperref}
\usepackage{tcolorbox}
\pagestyle{empty}
\setlength{\parindent}{0pt}
\setlength{\parskip}{3pt}

\begin{document}

% ---- Header --------------------------------------------------------------
\begin{center}
  {\LARGE\bfseries\color{aiNavy} AIinRT 2027}\\[2pt]
  {\small\color{aiGrey} Artificial Intelligence in Radiation Therapy
   \,\textperiodcentered\, 2027
   \,\textperiodcentered\, Princess M\'axima Center, Utrecht}\\[2pt]
  {\color{aiAccent}\bfseries Abstract submission template}\\[2pt]
  \textcolor{aiAccent}{\rule{\textwidth}{0.8pt}}
\end{center}

% ---- Instructions (DELETE THIS BOX BEFORE SUBMISSION) --------------------
\begin{tcolorbox}[colback=aiNavy!6,colframe=aiAccent,title=Instructions
  (delete this box before submission),fonttitle=\bfseries\small,fontupper=\small]
\begin{itemize}\setlength\itemsep{1pt}
  \item Main text max.\ \textbf{500 words}; \textbf{one A4 page} including
        display items and references.
  \item Max.\ \textbf{2 display items} (figures and/or tables combined);
        captions and references not counted in the word limit.
  \item Do not alter fonts, margins, or spacing; do not add pages.
  \item \textbf{Double-blind review}: no author names, affiliations,
        acknowledgements, or identifying information anywhere in this
        document; cite your own prior work in the third person. Author
        details are entered in the submission system only.
  \item \textbf{AI-use disclosure is mandatory}; ``None'' is acceptable.
  \item Submit as PDF via \url{https://AIinRT.org} (deadline: see website;
        2027 meeting dates to be confirmed). Accepted abstracts: proffered
        paper, 9\,min + 3\,min discussion.
\end{itemize}
\end{tcolorbox}

% ---- Title & authors ------------------------------------------------------
{\large\bfseries Title of the abstract (max.\ 20 words, sentence case, no
abbreviations)}\\[4pt]
{\small\itshape\color{aiGrey} Double-blind review --- no author names,
affiliations, or contact details in this document. Author information is
collected separately in the submission system.}

\section{Session preference (select one)}
{\small\itshape\color{aiGrey}
S1 Segmentation \& Registration \,\textperiodcentered\,
S2 Reconstruction \& Synthesis \,\textperiodcentered\,
S3 Foundation Models, Text, Explainability \& Uncertainty \,\textperiodcentered\,
S4 Dose \& Adaptive Workflows \,\textperiodcentered\,
S5 Clinical Predictions \& Outcomes \,\textperiodcentered\,
S6 Implementation, QA \& Ethics}

\section{Introduction / Purpose}
\textit{State the clinical or technical problem, its relevance to radiation
therapy, and the aim of the work. [Main text: max.\ 500 words in total across
the four sections.]}

\section{Materials \& Methods}
\textit{Describe data (cohort size, modality, source), model/algorithm,
training and validation strategy, and evaluation metrics. Report
ethics/consent status where applicable.}

\section{Results}
\textit{Present quantitative results with uncertainty where possible (e.g.,
mean\,$\pm$\,SD, CI). Refer to Figure~\ref{fig:example} / Table~1.}

\section{Conclusion}
\textit{Summarise the main finding and its potential clinical impact or next
steps.}

% ---- Display item (max. 2 in total) ---------------------------------------
\begin{figure}[h]
  \centering
  % \includegraphics[width=0.7\textwidth]{figure1}
  \fbox{\parbox[c][22mm][c]{0.7\textwidth}{\centering\itshape\color{aiGrey}
    Insert figure or table here (max.\ 2 display items in total)}}
  \caption{Concise, self-contained description (not counted in the 500 words).}
  \label{fig:example}
\end{figure}

\section{References (optional, max.\ 5)}
{\small [1] Author A \emph{et al.} Journal Year;Vol:Pages. DOI}

\section{Disclosure of AI use (mandatory)}
\textit{Declare any use of generative AI or AI-assisted tools in the conduct
of the research and/or the preparation of this abstract (e.g., code
generation, text drafting/editing, figure creation), naming the tool(s) and
purpose. Authors remain fully responsible for content. Write ``None'' if not
applicable.}

\section{Conflicts of interest \& funding}
\textit{Declare relevant conflicts of interest and funding sources, or state
``None''.}

\end{document}

